\documentclass[12pt]{article}
\usepackage[utf8]{inputenc}
\usepackage[T1]{fontenc}
\usepackage[margin=1in]{geometry}
\usepackage{setspace}
\usepackage{xcolor}
\usepackage[colorlinks=true, urlcolor=blue!60!black, linkcolor=blue!60!black]{hyperref}
\setlength{\parindent}{2em}
\setlength{\parskip}{1.25em}
\renewcommand{\baselinestretch}{1.0}
\begin{document}
\begin{center}
{\LARGE \textit{Rust v. Sullivan}, 500 U.S. 173 (1991)\par}
\vspace{0.8em}
{\large Case Brief --- Four-Page Edition\par}
\end{center}
\vspace{1em}

\section{Memory Jogger}

When the government appropriates public funds to establish a program promoting a particular policy, it may define the program's limits by prohibiting federally funded projects from engaging in abortion counseling, referral, or advocacy without violating the First or Fifth Amendments (\href{https://supreme.justia.com/cases/federal/us/500/173/}{500 U.S. 173}).

\section{Detailed Case Facts}

Title X of the Public Health Service Act (1970) funds voluntary family planning projects, but \S~1008 provides that ``[n]one of the funds appropriated under this subchapter shall be used in programs where abortion is a method of family planning.'' For about 18 years HHS read \S~1008 to bar performing abortions while permitting, and later requiring, nondirective counseling and referral.

In 1988, citing GAO and Inspector General reports, the Secretary issued regulations limiting Title X to preconceptional services and imposing three conditions. Projects could not counsel or refer for abortion as a method of family planning, even on request, and could respond that the project ``does not consider abortion an appropriate method of family planning''; they could not ``encourage, promote or advocate abortion''; and they had to be ``physically and financially separate'' from abortion activities.

\section{Procedural History}

Title X grantees and physicians, including Dr. Irving Rust, and the State and City of New York challenged the regulations on their face. The Southern District of New York upheld them, and a divided Second Circuit affirmed under \textit{Chevron}. The First and Tenth Circuits had invalidated them. The Supreme Court granted certiorari, heard argument October 30, 1990, and affirmed on May 23, 1991.

\section{Judicial Votes}

Affirmed, \textbf{5--4}.
\begin{itemize}
\item \textbf{Majority:} Chief Justice Rehnquist, joined by Justices White, Kennedy, Scalia, and Souter.
\item \textbf{Concurrences:} None.
\item \textbf{Dissents:} Justice Blackmun, joined by Justice Marshall, by Justice O'Connor as to Part I, and by Justice Stevens as to Parts II and III; Justice Stevens; Justice O'Connor.
\end{itemize}

\section{Holding}

Because \S~1008 is ambiguous, the regulations are a permissible construction entitled to \textit{Chevron} deference and do not raise ``grave and doubtful constitutional questions.'' The government may selectively fund a program and define its limits to exclude speech outside its scope without viewpoint discrimination or an unconstitutional condition, so long as the condition governs the funded \textit{project} rather than the \textit{grantee}. Refusing to subsidize abortion counseling places no obstacle in a woman's path beyond her indigency. Rehnquist cautioned that funding is not ``invariably sufficient'' to justify content control, flagging public forums and universities; the later ``government speech'' framing came in \textit{Rosenberger} and \textit{Velazquez}.

\section{Analysis of Opinions}

\textbf{Rehnquist (majority).} Facial challengers bear a heavy burden under \textit{Salerno}. The agency's reversal was supported by ``reasoned analysis.'' Under \textit{Maher}, \textit{Harris v. McRae}, and \textit{Regan}, the government may favor childbirth through selective funding, just as funding democracy abroad does not require funding rival philosophies. Unlike \textit{League of Women Voters}, the rules left grantees free to perform and advocate abortion through separate programs. They permitted emergency referrals, and doctors could say abortion advice is outside the program.

\textbf{Blackmun (dissenting).} The Court should have avoided the constitutional questions. The rules were viewpoint-based, compelled scripted statements, and intruded into truthful physician--patient dialogue; under \textit{Perry v. Sindermann} and \textit{Pickering}, workplace speech restrictions still require balancing. By manipulating the physician's words, the government distorted indigent women's choices under \textit{Roe}.

\textbf{Stevens (dissenting).} Read as a whole, Title X's ``PROHIBITION OF ABORTION'' targets conduct, not truthful speech, as 18 years of consistent agency construction confirmed.

\textbf{O'Connor (dissenting).} The rules raised ``serious First Amendment concerns,'' so the Court should have ruled on statutory grounds and left the question to Congress.

\section{Examples: Future Applications}

\begin{enumerate}
\item \textbf{Same-Side (Valid): Smoking-Cessation Grants.} Grant-funded counselors may promote nicotine replacement but not recommend vaping, while the clinic discusses vaping in separately funded programs. The government defines its own message.
\item \textbf{Same-Side (Valid): Democracy Curricula.} A civic-literacy grant requires curricula to advocate self-government over authoritarianism; recipients remain free to speak otherwise with other funds.
\item \textbf{Opposite-Side (Invalid): Organization-Wide Pledge.} Grantees must adopt a policy ``explicitly opposing prostitution and sex trafficking'' even with private funds. As in \textit{AID v. AOSI} (2013), this regulates speech outside the program.
\item \textbf{Opposite-Side (Invalid): Legal-Aid Restrictions.} Funded lawyers may not challenge welfare statutes' validity. As in \textit{Velazquez} (2001), the government funds private advocacy and distorts the courts.
\item \textbf{Fence-Sitter: University Research Grants.} Funded professors may not mention fossil-fuel phaseouts in grant-related work. Rehnquist flagged universities as ``traditional sphere[s] of free expression,'' so this sits on the contested frontier.
\end{enumerate}

\section{Critique}

\textbf{Conservative, libertarian, and textualist views.} Defenders stress democratic accountability: elected officials must control the message of programs taxpayers fund. Libertarians warn that, as the spending state grows, conditional grants let government buy silence; textualists note the statutory holding rested on \textit{Chevron} deference to a sharp policy reversal.

\textbf{Progressive and doctrinal views.} Dorothy Roberts, David Cole, and Robert Post criticize \textit{Rust} for enabling government control of knowledge and co-opting professionals to deliver state-scripted messages to low-income patients who expect candid advice. Others find the line between defining a program (\textit{Rust}) and imposing viewpoint conditions (\textit{Velazquez}, \textit{AOSI}) unstable.

\section{Key Quotations}

\begin{quote}
``The Government can, without violating the Constitution, selectively fund a program to encourage certain activities it believes to be in the public interest, without at the same time funding an alternate program which seeks to deal with the problem in another way. In so doing, the Government has not discriminated on the basis of viewpoint; it has merely chosen to fund one activity to the exclusion of the other.''
\end{quote}
\noindent --- Chief Justice Rehnquist, Opinion of the Court, 500 U.S. at 193.

\begin{quote}
``The Secretary's regulations do not force the Title X grantee to give up abortion-related speech; they merely require that the grantee keep such activities separate and distinct from Title X activities.''
\end{quote}
\noindent --- Chief Justice Rehnquist, Opinion of the Court, 500 U.S. at 196.

\begin{quote}
``When Congress established a National Endowment for Democracy to encourage other countries to adopt democratic principles, 22 U.S.C. \S~4411(b), it was not constitutionally required to fund a program to encourage competing lines of political philosophy such as Communism and Fascism.''
\end{quote}
\noindent --- Chief Justice Rehnquist, Opinion of the Court, 500 U.S. at 194.

\end{document}
